\section{Conclusión}

El presente hito completó la capa de presentación del sistema, cerrando el ciclo de 
desarrollo iniciado con el diseño de la base de datos y continuado con la 
implementación del backend. Las seis fases permitieron construir el frontend de 
forma incremental, donde cada módulo entregaba valor funcional inmediato y dejaba 
las condiciones necesarias para que el siguiente pudiera avanzar sin retrabajos.

Un aspecto que este hito deja en evidencia es la interdependencia real entre 
frontend y backend en un proyecto full-stack. La implementación de las vistas de 
productos y dispositivos fue la que reveló la ausencia de eliminaciones en cascada 
en la base de datos, forzando correcciones que debieron gestionarse mediante 
migraciones de Alembic. Este tipo de hallazgos, imposibles de anticipar durante el 
diseño inicial, reafirman la importancia de un proceso de desarrollo iterativo y de 
contar con herramientas como Alembic que permitan aplicar cambios controlados sin 
comprometer la integridad de los datos existentes.

Desde la perspectiva del usuario, el sistema resultante ofrece una experiencia 
diferenciada según el rol: el administrador dispone de un conjunto completo de 
herramientas para gestionar el inventario y los préstamos, mientras que el usuario 
común cuenta con una interfaz simplificada centrada en la consulta del catálogo 
disponible y el seguimiento de sus propias solicitudes. La fase de revisión de UX y 
UI garantizó que esta experiencia sea coherente tanto en escritorio como en 
dispositivos móviles, y que el sistema comunique de forma clara cada acción 
relevante a través de notificaciones y validaciones visuales.

Con este hito finalizado, el sistema cuenta con todas las capas del desarrollo 
completamente implementadas: modelo de datos, API REST y frontend. Los hitos 
siguientes abordarán la documentación del trabajo futuro, identificando las 
funcionalidades que quedan fuera del alcance del presente proyecto y describiendo 
cómo podrían llevarse a cabo.